% =========================================================================
% DRAFT — Appendix A: Muscle Force and Joint Torque Calculations
% Extracted from the actual MATLAB implementation on 2026-09-07.
% Every equation carries a \note{file:line} provenance comment so Ben can
% verify. This file is a DRAFT for review: paste sections into
% ProofFinal/chapters/92-AppendixA.tex (and Overleaf) only after approval.
% NOTE: appendix is currently \chapter{Appendix A: ...} per psutex style.
% =========================================================================

\chapter{Appendix A: Muscle Force and Joint Torque Calculations}\label{app:muscle_calcs}

This appendix documents the reference implementation of the force, torque, and
compliance equations used in Chapters~\ref{ch:methods}--\ref{ch:results}, as they
exist in the project's Matlab code. File names are given so that any number in this
dissertation can be traced to the code that produced it.

\section{BPA Force Model Implementation}

\subsection{Strain Measures}
The musculotendon path length $l_{LMT}$ is the sum of straight segments between
route via points, transformed into a common frame at each frame crossing. The
contractile strain used by the force model subtracts the artificial tendon and both
end-cap fittings from the path length before normalization by the resting length,
\begin{equation}
    \varepsilon \;=\; \frac{l_{0} - \left(l_{LMT} - l_{ten} - 2\,l_{fitting} - \chi_{0}\right)}{l_{0}},
    \label{eq:app_strain}
\end{equation}
where $\chi_{0}$ is the constant length offset of
Section~\ref{sec:improved_model} (zero in the uncorrected model).
[Provenance: \texttt{MonoPamDataExplicit.m} get.Contraction, lines 210--217;
\texttt{minimizeFlxPin.m} lines 150--164. End-cap fitting length
$l_{fitting}=\qty{25.4}{\mm}$ is a constructor default, line 69.]

The maximum free-load contraction at \qty{620}{\kPa} is
$\varepsilon_{620} = (l_{0}-l_{620})/l_{0}$, where $l_{620}$ is the BPA length at
\qty{620}{\kPa} with no external load (code variable \texttt{KMAX};
\texttt{MonoPamDataExplicit.m} line 278). The relative strain is
$\varepsilon^{*} = \varepsilon / \varepsilon_{620}$ (line 281).

\subsection{Normalized Force Surface}
The normalized isometric force surface of Equation~(\ref{eq:Fstar}) is implemented
in closed form in \texttt{Functions/f\_festo.m} (lines 1--26),
\begin{equation}
    F^{*} = c_{0}\left(e^{-c_{1}\varepsilon^{*}} - 1\right)
            + P^{*}\,e^{-c_{2}(\varepsilon^{*})^{2}},
    \qquad F^{*}\le 0 \;\Rightarrow\; 0,
    \label{eq:app_fstar}
\end{equation}
with $P^{*} = P/\qty{620}{\kPa}$. The coefficients were fitted with
\texttt{fittype} nonlinear least squares and least-absolute-residual robustness
(\texttt{Testing\_Data/bpaFits.m}, lines 45--60 and 115--140, lower bounds
$[0,0,0]$) and are reported in Table~\ref{tab:Fstar}. The code additionally holds a
40\,mm fit not used in this dissertation ($c_0 = 0.1224$, $c_1 = 10.47$,
$c_2 = 2.023$; \texttt{f\_festo.m} lines 17--20).

At runtime the torque model may alternatively evaluate a spline interpolation of
the manufacturer's datasheet curves (\texttt{Functions/festo4.m}, backed by
\texttt{FestoLookup.mat}), with force clamped to zero for $\varepsilon^{*}>1$ or
$F^{*}<0$. Equation~(\ref{eq:app_fstar}) replaced this lookup in the corrected
model.

\subsection{Maximum Force}
The resting-length-dependent maximum force of Equation~(\ref{eq:Fmax_3D}) is
implemented in \texttt{Functions/maxBPAforce.m} (10\,mm: line 46; 20\,mm: line 57;
duplicate for 20\,mm in \texttt{Functions/Fmax20.m}) with the fixed
\qty{7.5}{\mm} end-contact offset. The fitted 20\,mm coefficients carry confidence
intervals $a_1 \in [1.4745, 1.5011]$\,\si{\N\per\kPa} and
$a_2 \in [0.0238, 0.0258]$\,\si{\per\kPa\per\metre}
(\texttt{Fmax20.m} lines 26--27). Total force is the product
$F = F^{*}(\varepsilon^{*},P^{*})\,F_{620}(l_{0})$, resolved into a vector along the
force direction (\texttt{MonoPamDataExplicit.m} lines 283--287).

\section{Muscle Path Length and Moment Arm}

Path length accumulates the Euclidean norm of each segment between via points,
applying the body-frame transformation whenever the path crosses between the femur
and tibia frames (\texttt{MonoPamDataExplicit.m} lines 78--95, total at lines
122--133).

The vector moment arm is implemented as the perpendicular foot of the joint origin
onto the line of action,
\begin{equation}
    \vec{r} \;=\; \vec{p} \;-\; \hat{u}\left(\hat{u}\cdot\vec{p}\right),
    \label{eq:app_perpfoot}
\end{equation}
where $\vec{p}$ is the vector from the joint origin to the muscle-path crossing
point and $\hat{u}$ is the unit force direction
(\texttt{MonoPamDataExplicit.m} lines 156--168; identical implementations in
\texttt{MonoMuscleData.m} lines 111--121, \texttt{minimizeFlxPin.m} lines 400--409,
and \texttt{MonoPam\_mult.m} lines 186--199). This is the same quantity as the
cross-product projection of Equation~(\ref{eq:project_ma}); the planar scalar
moment arm about the joint axis is the magnitude of the in-plane components,
computed as \texttt{hypot(mA(:,1), mA(:,2))}
(\texttt{predictKneeFlexor20mm.m} line 61). Torque about the joint $z$-axis is
$\vec{M} = \vec{r}\times\vec{F}$, evaluated only where the BPA is in tension
(values are discarded for $\varepsilon < -0.02$, the compressive region;
\texttt{MonoPamDataExplicit.m} line 311).

\section{Compliance Model Implementation}

The compliance-corrected muscle length entering
Equation~(\ref{eq:Lm}) is computed as
$l_{m} = l_{LMT} - \chi_{0} - l_{ten} - 2\,l_{fitting}$
(\texttt{minimizeFlxPin.m} lines 254 and 368--377). The bracket stiffness matrix
and projected scalar compliance along the force direction are evaluated per
configuration as
$\mathbf{K}_{br} = \mathrm{diag}(\chi_{1},\chi_{2},\chi_{2})$ in the bracket frame
and $k_{\hat{u},br} = \left(\hat{u}^{\,T}\mathbf{K}_{br}^{-1}\hat{u}\right)^{-1}$
(\texttt{minimizeFlxPin.m} lines 266--272).
\textbf{Note for review:} the code sets the $z$-bracket element to $\chi_{2}$,
whereas Equation~(\ref{eq:bktstiffness}) in the text states
$\mathrm{diag}(\chi_{1},\chi_{2},\chi_{1})$; one of the two should be corrected
after checking the CAD frame convention (the extensor text in
Section~\ref{sec:improved_model} swaps the axes by design for the arch bracket).

The artificial tendon is modeled as an axial spring with stiffness
\begin{equation}
    k_{ten} \;=\; \frac{n\,A_{c}\,E}{L_{ten}},
    \label{eq:app_tendon}
\end{equation}
with $n=2$ parallel cable strands ($\times 6$ for the 20\,mm assembly),
$A_{c} = \qty{1.51e-6}{\metre\squared}$ (19-strand cable), and
$E = \qty{193e9}{\pascal}$
(\texttt{minimizeFlxPin.m} lines 460--471; \texttt{MonoPam\_mult.m} lines 606--617).
Series combination with the projected bracket compliance gives
$k_{eq} = \left(k_{\hat{u},br}^{-1} + k_{ten}^{-1}\right)^{-1}$
(lines 273--275).

The scalar compliance displacement $\delta$ of Equation~(\ref{eq:complianceforcebalance})
is found numerically with \texttt{fzero} on the bracket $[0,\, l_{m}-l_{620}]$,
\begin{equation}
    F_{620}(l_{0})\,F^{*}\!\left(\varepsilon^{*}(\delta),P^{*}\right) - k_{eq}\,\delta = 0,
    \label{eq:app_forcebalance}
\end{equation}
with $\delta \equiv 0$ when the uncorrected relative strain already exceeds unity
(\texttt{minimizeFlxPin.m} lines 291--305). The vector bracket deflection is then
$\boldsymbol{\epsilon}_{br} = \mathbf{K}_{br}^{-1}\left(\norm{\vec{F}}\hat{u}\right)$
(line 324), the cable stretch is $\norm{\vec{F}}/k_{ten}$ (line 334), and the
attachment point, path, moment arm, and torque are recomputed. For the two-BPA
flexor grouping, the coupled balance shares one common force between both muscles
and is solved analogously on $[0, F_{620}]$ (\texttt{MonoPam\_mult.m} lines
500--565).

\section{Wrapping-Loss Implementation}

For the pinned-knee extensor, Equation~(\ref{eq:deltaL}) is implemented
segment-wise over the wrap arc with radii $R_1 = \qty{12}{\mm}$ and
$R_2 = \qty{40}{\mm}$ and segment boundaries $\theta_{1} = \ang{27}$ and
$\theta_{2} = \ang{-23}$ (\texttt{minimizeExtX3.m}, nested \texttt{Contraction},
lines 327--370). For the biomimetic-knee 20\,mm cell, the arc-length factor is the
routed bend measure $b(\theta_{k})$ (arc length swept by the wrap, built with the
\texttt{wRap} bend radius in \texttt{buildKneeFlexorRoute20mm.m} lines 265--268),
giving $\Delta l = \chi_{3}\, b(\theta_{k})\,(1-\varepsilon^{*})^{2}$
(\texttt{MonoPam\_mult.m} lines 398--441), with a hardcoded-geometry fallback for
degenerate routes (lines 428--434). The loss enters only the force-producing
strain; the kinematic (prediction) strain excludes it
(\texttt{minimizeExtX3.m} lines 236--240).
